\section{根壁と anticanonical heart の古典的障害}
\label{sec:root-walls}
本節は第\ref{sec:six-point-classification}節の四交比トーラスの
境界を古典的な六点 blow-up で解釈する．任意の非可換 potential の
AS 性，あるいは Rains heart・VdB heart の境界特殊化を仮定しない．
\begin{proposition}[十一の divisor の根幾何]
\label{prop:root-dictionary}
$P=(Q^{\langle-1\rangle})^{\mathsf t}$ の列を $p_1,p_2,p_3$，
座標点を $p_4,p_5,p_6$ とする．十八係数がすべて非零のとき，
十一の禁止条件は次の有効根に対応する．
\begin{enumerate}
\item $Q$ の二行 $b,b'$ と二列 $c,c'$ による $2\times2$ minor は，
$\{0,1,2\}\setminus\{c,c'\}=\{j\}$ と書けば
$p_{b+1},p_{b'+1},p_{4+j}$ の共線性であり，
根は $h-e_{b+1}-e_{b'+1}-e_{4+j}$．全部で九つ．
\item $\det Q^{\langle-1\rangle}=0$ は $p_1,p_2,p_3$ の共線性で，
根は $h-e_1-e_2-e_3$．
\item $\det Q=0$ は六点が一つの conic 上にある条件で，
根は $2h-\sum_{i=1}^6e_i$．
\end{enumerate}
各根 $\alpha$ は $\alpha^2=-2$ と $H\cdot\alpha=0$
（$H=3h-\sum e_i$）を満たす．各 divisor の一般点の
真の変換は滑らかな有理 $(-2)$ 曲線 $E$ である．
\end{proposition}
\begin{proof}
二行二列の minor の消滅は，$P$ の二列と補集合の座標点を
並べた $3\times3$ 行列式の消滅と同値である．
$\det P=\det Q^{\langle-1\rangle}$ であり，
座標三点を通る conic の三可動点での評価行列は，
非零因子を除いて $Q$ に一致する（定理\ref{thm:general-position-iff}）．
交点数は $h^2=1,e_i^2=-1$ から従う．
一般点では余分な共線性・重複点・conic の退化はない．
\end{proof}

\begin{theorem}[一つの根壁では導来圏も異なる]
\label{thm:classical-root-wall}
十一因子の一般点に対応する六点 blow-up を
$S=\operatorname{Bl}_{p_1,\ldots,p_6}\mathbb P^2$，
対応する $(-2)$ 曲線を $i:E\hookrightarrow S$ とする．
$H=-K_S$ と置く．
\begin{enumerate}
\item $S$ は滑らかな弱 del Pezzo 曲面で $H^2=3$，
$H$ は nef かつ big だが ample ではない．反標準射
\[
f:S\longrightarrow X=\operatorname{Proj}\bigoplus_{n\ge0}H^0(S,nH)
\]
は $E$ を収縮し，$X$ は $A_1$ 特異点を持つ正規三次曲面である．
\item $F=i_*\mathcal O_E(-1)$ は非零な $2$-spherical object であり，
\[
F\otimes\mathcal O_S(H)\simeq F,\qquad
\operatorname{RHom}_S(\mathcal O_S(nH),F)=0\quad(n\in\mathbb Z),
\qquad \mathbf Rf_*F=0.
\]
従って $(\mathcal O_S,(-)\otimes\mathcal O_S(H))$ は
$\operatorname{coh}S$ において Artin--Zhang ample でない．
\item $R_S=\bigoplus_{n\ge0}H^0(S,nH)$ と置くと
$\operatorname{qgr}R_S\simeq\operatorname{coh}X$ だが，
\[
\boxed{D^b(\operatorname{coh}S)\not\simeq
       D^b(\operatorname{qgr}R_S).}
\]
従ってこの古典的境界は，導来同値を前提とする
単なる HRS tilt の例ではない．
\end{enumerate}
\end{theorem}
\begin{proof}
十一因子の一般点では六点は相異なり，四点共線などの
追加退化はない．六点 blow-up に対する弱 del Pezzo の判定と，
反標準射が $(-2)$ 曲線を rational double point に収縮する定理を使う
\cite[Introduction, \S4]{BlancLamy}
\cite[Theorem 8.3.2]{Dolgachev}．
他の有効根のない一般点なので特異点は $A_1$ 型である．
$H\cdot E=0$ から $H$ は ample ではない．

$E\simeq\mathbb P^1$ かつ $E^2=-2$ より
$N_{E/S}\simeq\mathcal O_E(-2)$ である．
局所 Ext 層 $\mathcal Ext_S^0(F,F)=\mathcal O_E$，
$\mathcal Ext_S^1(F,F)=N_{E/S}$ とその局所大域スペクトル列から
\[
\operatorname{Ext}^j_S(F,F)=
\begin{cases}k&j=0,2,\\0&j\ne0,2.\end{cases}
\]
$K_S|_E\simeq\mathcal O_E$ なので Serre 関手は
$F\mapsto F[2]$ となり，$F$ は $2$-spherical．
同じ理由で $H|_E\simeq\mathcal O_E$ だから
\[
\operatorname{RHom}_S(\mathcal O_S(nH),F)
\simeq\mathbf R\Gamma(E,\mathcal O_E(-1))=0
\]
がすべての $n$ で成り立つ．
非零対象 $F$ が全反標準 twist から見えないので
AZ の生成条件は失敗する．
$f(E)$ は一点で $H^0(E,\mathcal O_E(-1))
=H^1(E,\mathcal O_E(-1))=0$ より $\mathbf Rf_*F=0$．

$X$ は $A_1$ 特異点を持つ正規三次曲面であり，
その homogeneous coordinate ring の $\operatorname{qgr}$ は
$\operatorname{coh}X$ である．滑らかな $S$ では
有界 coherent 複体の任意の二対象間の高次 Ext は消滅する．
ところが $X$ の特異点 $x$ の非正則局所環では
剰余体 $k(x)$ の最小自由解消が無限に続き，
$\operatorname{Ext}_X^j(k(x),k(x))\ne0$ となる $j$ が
無限に存在する．この性質は三角圏の同値と両立しない．
\end{proof}

\begin{corollary}[任意次数の弱 del Pezzo に対する障害]
\label{cor:all-degree-weak-dp}
標数零の代数閉体上の任意次数 $d=K_S^2\ge1$ の滑らかな
弱 del Pezzo 曲面 $S$ に $(-2)$ 曲線 $E$ があるとする．
$\tau=(-)\otimes\mathcal O_S(-K_S)$ に対し
$(\mathcal O_S,\tau)$ は AZ ample ではなく，
\[
D^b(\operatorname{coh}S)\not\simeq
D^b\!\left(\operatorname{qgr}
 \bigoplus_{n\ge0}H^0(S,-nK_S)\right).
\]
\end{corollary}
\begin{proof}
任意次数の弱 del Pezzo の反標準モデルは $(-2)$ 曲線を
Du Val 特異点へ収縮する\cite[Theorem 8.3.2]{Dolgachev}．
同じ $F=i_*\mathcal O_E(-1)$ の全反標準 twist からの
導来 Hom が零となること，および反標準モデルの特異点で
剰余体が非有界の Ext を持つことを使う．
\end{proof}

\begin{proposition}[一つの根壁の有理数上の明示例]
\label{prop:explicit-single-root-wall}
\[
Q_0=\begin{pmatrix}1&1&1\\1&2&2\\1&3&4\end{pmatrix}
\]
では行 $(0,1)$・列 $(1,2)$ の minor だけが零であり，
$\det Q_0=1$，$\det Q_0^{\langle-1\rangle}=1/24$．
六点は
\[
(1:1:1),(2:1:1),(12:4:3),(1:0:0),(0:1:0),(0:0:1).
\]
共線三点は $p_1,p_2,p_4$ だけで，直線は $y=z$．
$|-K_S|$ の基底を
\[
\begin{aligned}
 A&=\tfrac13x^2(y-z)-xy(y-z),\\
 B&=\tfrac7{12}x^2(y-z)-x(y^2-z^2),\\
 C&=\tfrac16x^2(y+2z)-\tfrac32xy^2+y^2z,\\
 D&=\tfrac14x^2(y+z)-\tfrac32xy^2+yz^2
\end{aligned}
\]
とすると反標準像の三次方程式は
\[
\begin{aligned}
0={}&23A^3-22A^2B+44A^2C-42A^2D+8AB^2-16ABC\\
&-24AC^2+24AD^2+8B^2C+24BC^2-24BCD .
\end{aligned}
\]
曲線 $y=z$ の像は $[A:B:C:D]=[0:0:1:1]$．
この点で勾配は零であり，$C=1$ の局所座標
$(A,B,D-1)$ の Hessian 行列式は $-2304\ne0$．
従って実際に通常二重点である．
\end{proposition}

\begin{remark}[非可換への適用範囲]\label{rem:noncommutative-wall-scope}
本節は古典的 blow-up の coherent heart とその
anticanonical heart の比較である．十一因子上の一般の非可換
potential が AS heart を定めるか，Rains heart・VdB heart の
特殊化が $\operatorname{coh}S$ と同定できるかは未証明．
特に VdB の一般位置での導来同値をこの根壁まで
無条件に延長してはならない．本例では
$\mathbf Rf_*$ 自身が非零対象を零に送るので，
単なる振幅一の同値に伴う HRS tilt とは根本的に異なる．
\end{remark}
